
\subsubsection{5.1 RQ1: H_d Signal Discriminability (H-E1)}

\textbf{Finding: Domain-specific H_d signals discriminate saturated from healthy benchmarks with very large effect sizes on synthetic panels calibrated to real PWC statistics. Real-data validation has not been performed.}

\textbf{Table 1: H-E1 Signal Discriminability Results (Synthetic Panels)}

\begin{table}[htbp]
\centering
\begin{tabular}{lllllll}
\toprule
Domain & Mann-Whitney p & Cohen's \ & d\ &  & AUC (standard direction) & Gate \\
\midrule
CV & < 0.0001 & 5.267 & 0.000 (direction inverted; corrected AUC = 1.000) & PASS &  &  \\
NLP & < 0.0001 & 6.910 & 1.000 & PASS &  &  \\
Tabular & < 0.0001 & 6.515 & 1.000 & PASS &  &  \\
\bottomrule
\end{tabular}
\end{table}

\textit{Gate criterion: p < 0.05 AND |d| > 0.5 in at least 2/3 domains. All three domains pass.}

\textit{All results are from synthetic benchmark panels (20 saturated + 20 healthy per domain) calibrated to real PWC statistics. Effect sizes on real benchmark data have not been measured and may differ.}

The CV AUC appears as 0.000 in the standard direction because CV H_d (score variance) is \textit{lower} for saturated benchmarks, opposite to NLP and tabular. When the direction is corrected, the discriminability is perfect (AUC = 1.000), consistent with the |d| = 5.267 effect size. This directional asymmetry is discussed further in Section 5.5.

\begin{figure}[htbp]
\centering
\includegraphics[width=\columnwidth]{figures/boxplots.png}
\caption{H_d signal discriminability (boxplots, synthetic panels)}
\label{fig:boxplots}
\end{figure}

\textbf{Effect sizes increase monotonically with lookback window length.} Effect sizes were computed at 6, 12, 18, and 24 months of lookback:

\textbf{Table 2: Cohen's |d| by Lookback Window (Synthetic Panels)}

\begin{table}[htbp]
\centering
\begin{tabular}{llllllllll}
\toprule
Lookback & CV \ & d\ &  & NLP \ & d\ &  & Tabular \ & d\ &  \\
\midrule
6 months & 2.50 & 3.28 & 3.10 &  &  &  &  &  &  \\
12 months & 3.42 & 4.49 & 4.24 &  &  &  &  &  &  \\
18 months & 4.35 & 5.70 & 5.38 &  &  &  &  &  &  \\
24 months & 5.27 & 6.91 & 6.52 &  &  &  &  &  &  \\
\bottomrule
\end{tabular}
\end{table}

Effect sizes increase monotonically in all three domains across lookback window lengths, supporting the use of a 24-month lookback window as the maximal discriminability configuration.

\begin{figure}[htbp]
\centering
\includegraphics[width=\columnwidth]{figures/temporal_separation.png}
\caption{Effect size vs. lookback window}
\label{fig:temporal_separation}
\end{figure}

\subsubsection{5.2 RQ2: Granger-Predictive Mechanism — Submissions → Compression (H-M1)}

\textbf{Finding: Submission accumulation Granger-causes score variance compression at lag 2 (approximately 6 months; p = 1.854 × 10⁻⁵), with reverse causality not confirmed. 31.1\% of qualifying benchmarks exhibit compression.}

\textbf{Table 3: H-M1 Mechanism Validation Results}

\begin{table}[htbp]
\centering
\begin{tabular}{llll}
\toprule
Metric & Value & Target & Status \\
\midrule
Raw archive rows & 48,311 & — & — \\
Panel observations & 6,938 & — & — \\
Panel benchmarks & 466 & — & — \\
σ_measurement (median, 7,592 benchmarks) & 0.3323 & — & — \\
Compression events detected & 389 & — & — \\
Benchmarks with ≥ 1 compression event & 145 / 466 & — & 31.1\% \\
Spearman ρ (cumulative_count vs. score_var_top10) & 0.052 & > 0.4 & Below target \\
Spearman p & 1.51 × 10⁻⁵ & < 0.05 & Significant \\
Benchmarks tested for Granger (≥ 9 quarters) & 41 & ≥ 30 & Met \\
Granger p at lag = 2 (minimum across 41 benchmarks) & **1.854 × 10⁻⁵** & < 0.05 & **PASS** \\
Reverse Granger p at lag = 2 & > 0.05 & — & Not confirmed \\
Benchmarks significant individually (Granger, lag = 2) & 5 / 41 (12.2\%) & — & — \\
\bottomrule
\end{tabular}
\end{table}

The gate logic for H-M1 was defined as: Spearman (ρ > 0.4 AND p < 0.05) OR Granger (p < 0.05 at lag = 2). The Spearman path failed (ρ = 0.052 < 0.4); the Granger path passed (p = 1.854 × 10⁻⁵ < 0.05). The gate result is PASS via the Granger path.

The minimum p = 1.854 × 10⁻⁵ across 41 independent Granger tests passes Bonferroni correction for 41 simultaneous tests (corrected threshold: 0.05/41 ≈ 0.00122).

The dissociation between low linear correlation (Spearman ρ = 0.052) and strong Granger signal (p = 1.854 × 10⁻⁵) is consistent with a threshold-triggered nonlinear mechanism: compression events occur when cumulative submission count crosses a benchmark-specific saturation threshold, producing a weak contemporaneous correlation but a detectable lagged structure. Only 12.2\% of individually tested benchmarks show significant Granger results; this rate is likely a lower bound, as many benchmarks were excluded from individual testing due to insufficient time series length (< 9 quarters).

\begin{figure}[htbp]
\centering
\includegraphics[width=\columnwidth]{figures/lag_profile.png}
\caption{Granger p-value profile across lags 1–4, forward and reverse directions}
\label{fig:lag_profile}
\end{figure}

\begin{figure}[htbp]
\centering
\includegraphics[width=\columnwidth]{figures/compression_distribution.png}
\caption{Distribution of compression events per benchmark}
\label{fig:compression_distribution}
\end{figure}

\begin{figure}[htbp]
\centering
\includegraphics[width=\columnwidth]{figures/reverse_causality.png}
\caption{Forward vs. reverse Granger p-value comparison}
\label{fig:reverse_causality}
\end{figure}

\subsubsection{5.3 RQ3: Cross-Sectional Diagnostic Utility (H-M2)}

\textbf{Finding: NLP H_d achieves AUC = 0.857 for cross-sectional discrimination of compressed from non-compressed benchmarks on real data. CV and tabular do not meet the AUC > 0.8 gate. The temporal ordering gate failed (Section 5.4).}

\textbf{Table 4: H-M2 Cross-Sectional Diagnostic Results (Real Data)}

\begin{table}[htbp]
\centering
\begin{tabular}{lllll}
\toprule
Domain & Mann-Whitney p & AUC_lead & AUC_concurrent & Gate (AUC > 0.8) \\
\midrule
CV & 1.000 & 0.390 & 0.564 & Fail \\
NLP & **0.0076** & **0.857** & 0.835 & Pass \\
Tabular & 0.0435 & 0.318 & 0.318 & Fail \\
\bottomrule
\end{tabular}
\end{table}

NLP H_d discriminates compressed from non-compressed benchmarks with AUC = 0.857 on real panel data. CV and tabular do not achieve the AUC > 0.8 gate. These are cross-sectional results and do not establish temporal precedence (see Section 5.4).

\begin{figure}[htbp]
\centering
\includegraphics[width=\columnwidth]{figures/auc_comparison.png}
\caption{AUC comparison: lead vs. concurrent, per domain}
\label{fig:auc_comparison}
\end{figure}

\begin{figure}[htbp]
\centering
\includegraphics[width=\columnwidth]{figures/mann_whitney_boxplot.png}
\caption{Mann-Whitney boxplot: H_d magnitude, compressed vs. non-compressed}
\label{fig:mann_whitney_boxplot}
\end{figure}

\subsubsection{5.4 Temporal Ordering: Operationalization Failure (H-M2 Gate)}

\textbf{Finding: The H-M2 temporal ordering gate failed. Only 1 collapse event was detected under the pre-registered expanding Kendall τ criterion (minimum required: 20). All three domains produced fraction_leading = 0.000. This result reflects an operationalization failure, not a conclusion about the underlying mechanism.}

\textbf{Table 5: H-M2 Primary Gate Results}

\begin{table}[htbp]
\centering
\begin{tabular}{llll}
\toprule
Domain & fraction_leading & Collapse events & Gate (≥ 0.60) \\
\midrule
CV & 0.000 & 0 & Fail \\
NLP & 0.000 & 0 & Fail \\
Tabular & 0.000 & 0 & Fail \\
**Overall** & — & **1 total** & **Fail** \\
\bottomrule
\end{tabular}
\end{table}

The gate required fraction_leading ≥ 0.60 in at least 2 domains. All domains produced fraction_leading = 0.000 because only 1 collapse event was detected even after a post-hoc mitigation that lowered the τ threshold from 0.90 to 0.85.

\textbf{Root cause.} The expanding Kendall τ collapse criterion requires \texttt{score_var_top10} to increase monotonically over time. However, H-M1 established that \texttt{score_var_top10} \textit{decreases} for compressed benchmarks (compression is the reduction in variance). These two signals are anti-correlated by construction: the same column cannot simultaneously fall (as it does during compression) and rise monotonically (as required by the expanding τ criterion). Ablation across all five variants (A1–A5) produced 0 onset events; only when the pool was expanded to all benchmarks without compression filtering did 1 event appear.

This is a structural incompatibility between the operationalization and the data, not evidence that H_d signals fail to precede collapse. The cross-sectional diagnostic results (NLP AUC = 0.857, Table 4) remain valid, but temporal precedence was not established.

\textbf{H-M3 and H-M4 were not executed} as a consequence of the H-M2 gate failure. Building a Cox PH survival model on a dataset with 1 collapse event would not produce interpretable results.

\begin{figure}[htbp]
\centering
\includegraphics[width=\columnwidth]{figures/km_lead_time_curves.png}
\caption{Kaplan-Meier survival curves (fully censored due to insufficient collapse events)}
\label{fig:km_lead_time_curves}
\end{figure}

\subsubsection{5.5 Domain Saturation Phenotype Asymmetry}

\textbf{Finding: CV saturation manifests as score convergence (variance decreases), while NLP and tabular saturation manifests as signal divergence (deviation or rank correlation increases). This directional asymmetry was not anticipated in the original hypothesis design and requires per-domain sign normalization for any multi-domain combined score.}

This asymmetry was observed empirically in H-E1 (CV AUC = 0.000 in standard direction, = 1.000 when direction-corrected) and confirmed in H-M1 (compression events defined by variance decrease). Two interpretations are consistent with the data:

\begin{itemize}
\item \textbf{Convergence phenotype (CV):} Models overfit the test distribution, producing increasingly similar scores; variance falls as the score distribution compresses around a ceiling.
\item \textbf{Divergence phenotype (NLP/tabular):} Contamination or protocol gaming produces anomalously high scores, increasing deviation from the bulk distribution.
\end{itemize}

The implication for the BCBHS Cox model design is that a naive combined covariate would aggregate incoherent directional signals. Per-domain sign normalization is required before multi-domain combination.

\subsubsection{5.6 Summary Evidence Table}

\textbf{Table 6: Status of All Claims}

\begin{table}[htbp]
\centering
\begin{tabular}{lllllll}
\toprule
Claim & Evidence & Status &  &  &  &  \\
\midrule
H_d signals discriminate ( & d & > 0.5, p < 0.05) &  & d & > 5 all 3 domains, synthetic panels (H-E1) & Confirmed (synthetic data only) \\
H_d signals discriminate on real data & Not measured via H_d on real compression labels & Unverified &  &  &  &  \\
Submission accumulation Granger-causes compression & Granger p = 1.854 × 10⁻⁵, lag = 2 (H-M1) & Confirmed &  &  &  &  \\
31.1\% benchmark compression prevalence & 389 events, 145/466 benchmarks (H-M1) & Confirmed &  &  &  &  \\
NLP cross-sectional AUC > 0.8 & AUC = 0.857 (H-M2) & Confirmed &  &  &  &  \\
H_d signals temporally precede collapse (≥ 12 months) & fraction_leading = 0.0 all domains; 1 event (H-M2) & Refuted by operationalization failure; mechanism not tested &  &  &  &  \\
Cox PH C-index ≥ 0.70 & H-M3 not executed & Inconclusive &  &  &  &  \\
Lead time ≥ 12 months (Kaplan-Meier) & H-M4 not executed & Inconclusive &  &  &  &  \\
\bottomrule
\end{tabular}
\end{table}

